%% 03-key-findings.tex — Five headline numbers

\section{Five Headline Numbers}

The portfolio distills to five numbers, each the headline result of
a specific paper or cluster of papers.

\begin{description}[leftmargin=0pt, itemsep=10pt]

  \item[\Large +22\% compactness improvement.]
    Across all 435 congressional districts, algorithmic maps are
    22\% more compact than the enacted 2020 maps, measured by the
    Polsby-Popper score (ratio of district area to perimeter squared).
    Compact districts follow natural geographic features rather than
    being stretched to capture partisan voters.
    \hfill\textit{Paper B.2}~\cite{dellaLibera2026edgeWeighted}

  \item[\Large North Carolina: 7 Democratic, 7 Republican.]
    When the algorithm partitions North Carolina---a state that has
    been the subject of repeated gerrymandering litigation---it
    produces a 7-to-7 delegation, closely matching the state's
    roughly even partisan split.  No partisan data was used.
    \hfill\textit{Paper B.11}~\cite{dellaLibera2026ncRegions}

    \smallskip
    \noindent\textit{Note: This result holds under the ApportionRegions algorithm
    ($k=14=7\times2$); the same data under GeoSection gives 5D/9R, illustrating
    that algorithm choice---not the data---determines partisan outcomes.}

  \item[\Large $T = 600$ bisection steps to convergence.]
    The algorithm reliably converges---produces stable, repeatable
    districts---within 600 METIS bisection iterations.
    More iterations change nothing; fewer can leave the solution
    incomplete.  This gives practitioners a firm stopping rule.
    \hfill\textit{Paper B.16}~\cite{dellaLibera2026convergence}

  \item[\Large 42\% minority-population threshold for VRA compliance.]
    States where minority voters make up at least 42\% of the
    voting-age population can reliably achieve proportional
    majority-minority representation through algorithmic redistricting.
    Below this threshold, geography alone makes proportional
    representation infeasible---regardless of who draws the map.
    \hfill\textit{Paper D.1}~\cite{dellaLibera2026vraCompliance}

    \smallskip
    \noindent\textit{This is an empirical regularity derived from 43-state analysis,
    not a legal bright line; VRA Section~2 compliance still requires the full
    Gingles three-prong test.  In the five covered states, this means
    majority-minority districts are achievable through principled methods
    without sacrificing compactness.}

  \item[\Large $O(n^{1.07})$ runtime.]
    The algorithm scales nearly linearly with the number of census
    tracts.  Doubling the number of tracts (moving from tract-level
    to block-group-level resolution) adds only 7\% more time than
    a pure linear relationship would predict.  All 50 states run
    in under eighteen minutes on commodity hardware.
    \hfill\textit{Paper B.6}~\cite{dellaLibera2026complexity}

\end{description}

\bigskip
\noindent These five numbers are conservative: they are derived from
runs on publicly available Census data using openly published code.
Any researcher with an internet connection can reproduce them.
